% ==========================================================
% titles/title-07.tex -- Title VII: External Relations
% ==========================================================

\ConstTitle{External Relations}{title:external-relations}

\Article{The Foreign Relations Council}{art:foreign-relations}

\ArtSection{Composition}{art:foreign-relations:s-composition}
The \gls{gls:foreign-relations-council} is a standard body of twelve members serving three-year terms (staggered so that four rotate out annually), selected from a qualification pool with demonstrated relevant expertise in international law, economics, diplomacy, regional affairs, and languages.

\ArtSection{Function}{art:foreign-relations:s-function}
The Council manages ongoing diplomatic relationships, negotiates proposed international agreements, and represents the Commonwealth in bilateral foreign relations. All negotiations are logged in real time by the \gls{gls:transparency-engine}. Sensitive ongoing negotiations may be withheld from public disclosure for a maximum of one year. In multilateral settings where the \gls{gls:residual-coordination-council} provides Commonwealth representation under Article \ref{art:residual-coordination}, the \gls{gls:foreign-relations-council} retains primacy over all substantive positions articulated by that representation. The Council must provide timely sign-off on substantive positions to prevent diplomatic delay; unreasonable delay in sign-off may be referred to the \gls{gls:constitutional-court} for expedited review.

\ArtSection{Mandate Constraints}{art:foreign-relations:s-mandate-constraints}
The Council may negotiate but not conclude binding international agreements. All treaties require ratification by the \gls{gls:sortition-legislature}. The Council is explicitly prohibited from negotiating any agreement that creates dependency relationships or coercive power over another nation's survival or self-determination. It has no authority over domestic policy under any circumstances.

\ArtSection{Conflict of Interest}{art:foreign-relations:s-conflict-interest}
Members of the Council are subject to the same conflict of interest, disclosure, and sector exclusion provisions applicable to all governance roles under Article \ref{art:governance-principles}, \ref{art:governance-principles:s-conflict-interest-sector}, and the same post-service employment prohibitions as \gls{gls:sortition-legislature} members under Article \ref{art:sortition-legislature}, \ref{art:sortition-legislature:s-post-service-employment}. Members may not hold financial interests in any entity with active commercial relationships with nations under active negotiation.

\Article{The Anti-Imperialism Principle}{art:anti-imperialism}

\ArtSection{Constitutional Prohibition}{art:anti-imperialism:s-constitutional-prohibition}
The Commonwealth may not enter into any relationship with any other nation, people, or international institution that creates dependency, extracts value without equivalent return, or gives the Commonwealth coercive power over the survival or self-determination of others.

\ArtSection{Specific Prohibitions}{art:anti-imperialism:s-specific-prohibitions}
\begin{itemize}
  \item Maintaining military bases or permanent military presence in other countries
  \item Using currency, credit, or financial system access as instruments of foreign policy coercion
  \item Holding other nations' debt as leverage
  \item Providing conditional aid that requires policy compliance from recipients
  \item Conducting covert operations that interfere in other nations' internal governance
  \item Entering trade relationships designed to create dependency in the other party
\end{itemize}

\ArtSection{Tariff Policy}{art:anti-imperialism:s-tariff-policy}
Ordinary protective or revenue tariffs are not prohibited by this Article and are distinct from the dependency-creating and coercive trade practices barred under \ref{art:anti-imperialism:s-specific-prohibitions}. Any tariff must be narrow in scope, tied to a specific and publicly stated purpose -- such as protecting Worker Cooperative Sector wages and labor standards under Article \ref{art:worker-coop-sector} from underpriced production that would not meet those standards domestically, or shielding a defined nascent domestic industry for a stated developmental purpose -- and must include a sunset date at enactment, not to exceed ten years, renewable only by an affirmative Legislative vote with published justification. A tariff may not be maintained indefinitely by repeated renewal absent a renewed and specific finding of ongoing purpose. A tariff structured or applied to create dependency in a trading partner, to coerce policy compliance, or to substitute for the prohibitions in \ref{art:anti-imperialism:s-specific-prohibitions} is prohibited regardless of its stated purpose.

\ArtSection{Positive Obligations}{art:anti-imperialism:s-positive-obligations}
The Commonwealth actively shares scientific knowledge, medical research, agricultural technology, and renewable energy systems as global Commons. Humanitarian aid is provided whenever possible and without conditions. The Commonwealth participates in and advocates for genuinely multilateral international institutions with equal voice regardless of economic size. Indigenous nations with whom the Commonwealth has completed Land Back agreements under Article \ref{art:indigenous-sov} are among the priority relationships under these positive obligations; sovereignty recognition is the beginning of a relationship, not the end of one.

\ArtSection{Collective Defense}{art:anti-imperialism:s-collective-defense}
The Commonwealth may participate in genuinely multilateral defense arrangements provided they do not require the development of offensive capability prohibited under Article \ref{art:defense}, the establishment of permanent military presence outside Commonwealth territory, or the subordination of Commonwealth foreign policy to any other nation's direction. The tension between collective defense obligations and the anti-imperialism principle is acknowledged; the \gls{gls:treaty-review-body} under \ref{art:anti-imperialism:s-treaty-review} assesses all proposed collective defense arrangements for compliance with both.

\ArtSection{Treaty Review}{art:anti-imperialism:s-treaty-review}
An independent \gls{gls:treaty-review-body}, a standard body of twenty members serving four-year terms selected from the international law expert qualification pool, reviews all proposed international agreements before ratification for compliance with the anti-imperialism principle and \ref{art:anti-imperialism:s-collective-defense}. Any agreement found to create dependency relationships or coercive power over another nation may not be ratified regardless of legislative majority.

\ArtSection{Human Rights and Non-Interference}{art:anti-imperialism:s-human-rights-non}
The Commonwealth acknowledges that the anti-domination principle does not stop at its borders. Where persons in neighboring or other states experience systematic domination by their own governments or other powerful actors, the Commonwealth's commitment to human rights generates obligations of solidarity that must be reconciled with the anti-imperialism principle of this Article.

The Commonwealth fulfills these obligations exclusively through non-coercive means: unconditional welcome to persons fleeing persecution under Article \ref{art:open-presence}; provision of humanitarian aid, medical assistance, and knowledge sharing without conditions under \ref{art:anti-imperialism:s-positive-obligations}; advocacy within genuinely multilateral institutions for universal human rights standards; and material support to non-governmental mutual aid organizations and human rights documentation bodies operating in affected nations.

The Commonwealth may not fulfill solidarity obligations through economic sanctions, conditional aid requiring governance changes, funding of political opposition or regime change organizations, covert operations, or military action regardless of humanitarian framing. These prohibitions apply without exception. The use of humanitarian language to justify interference is itself a constitutional violation.

The \gls{gls:foreign-relations-council} (Article \ref{art:foreign-relations}) determines whether any proposed solidarity action constitutes permissible non-coercive support or prohibited interference, against published criteria developed in consultation with the \gls{gls:treaty-review-body}. Determinations are subject to \gls{gls:constitutional-court} review. All solidarity activities and their funding are logged in full through the \gls{gls:transparency-engine}.

The Commonwealth maintains transparent public documentation of human rights conditions in all nations with which it has diplomatic relationships, published annually through the \gls{gls:transparency-engine}. This documentation serves the Commonwealth's civic membership and the international community; it is not an instrument of foreign policy pressure.

\Article{Defense and Nuclear Posture}{art:defense}

\ArtSection{Defensive Sufficiency}{art:defense:s-defensive-sufficiency}
The Commonwealth maintains military capacity sufficient for genuine territorial defense and no more. Force design is oriented exclusively toward defense. The Commonwealth does not maintain offensive power projection capacity beyond its borders.

\ArtSection{Doctrinal Orientation}{art:defense:s-doctrinal-orientation}
Force design is explicitly structured around denial rather than punishment. Ground forces are optimized for territorial defense, not expeditionary deployment. Air forces are oriented toward air superiority over Commonwealth territory, not long-range strike capability. Naval forces are configured for coastal and exclusive economic zone defense, not blue-water power projection. Long-range precision strike systems whose primary design purpose is offensive are prohibited.

\ArtSection{Military Governance}{art:defense:s-military-governance}
Military leadership is civilian-accountable. The \gls{gls:implementation-councils} overseeing defense (Article \ref{art:implementation-councils}) is civilian, selected by sortition, with full access to all military operations and decisions. The military budget is constitutionally capped as a percentage of GDP under \ref{art:defense:s-military-budget-cap} and may not be increased by simple legislative majority.

\ArtSection{Nuclear Posture}{art:defense:s-nuclear-posture}
The Commonwealth possesses the minimum number of nuclear weapons necessary to maintain credible second-strike capability. The arsenal is constitutionally capped and designed exclusively for second-strike purposes. First-strike capable systems are prohibited.

\ArtSection{No-First-Use}{art:defense:s-no-first-use}
The Commonwealth will never use nuclear weapons against a non-nuclear state under any circumstances. The Commonwealth will never use nuclear weapons first against a nuclear-armed state. This commitment is constitutionally entrenched and cannot be overridden by any emergency declaration, executive decision, or simple legislative majority.

\ArtSection{Nuclear Use Authorization Council}{art:defense:s-nuclear-use-authorization}
A \gls{gls:nuclear-use-authorization-council} of eight persons, selected by sortition and serving staggered two-year non-renewable terms in continuous session with no other governance role, holds sole nuclear launch authorization. Authorization thresholds:
\begin{itemize}
  \item Verified incoming nuclear attack confirmed by independent sensors: six of eight members required
  \item Any ambiguous or escalatory scenario: unanimous authorization required plus \gls{gls:constitutional-court} emergency review, completed within the time the Council determines remains before authorization would otherwise lapse under \ref{art:defense:s-authorization-window}; if review cannot be completed in that time, authorization does not issue
  \item Any first use scenario: constitutionally prohibited regardless of vote
\end{itemize}

\ArtSection{The Authorization Window}{art:defense:s-authorization-window}
Unanimous authorization under an ambiguous or escalatory scenario lapses automatically if not acted upon, and does not carry over to a materially changed situation. As with the seventy-two-hour gap in Article \ref{art:emergency-protocols}, this Constitution acknowledges that compressing constitutional review into a live nuclear crisis is a genuine vulnerability that cannot be fully designed away. The fail-safe default in every case where required review is incomplete, doubtful, or contested is that authorization does not issue.

All eight members are publicly identified, subject to full transparency requirements, and recallable by legislative supermajority.

\ArtSection{Defensive Posture Review}{art:defense:s-defensive-posture-review}
The \gls{gls:defensive-posture-review-body}, a standard body of sixteen members serving four-year terms selected from the strategic and arms control studies qualification pool, conducts mandatory public review of all military force structure, doctrine, procurement, and capability every four years. The Body classifies all major capabilities as defensive-primary or dual-use through a published Capability Classification Protocol; offensive-primary systems are constitutionally prohibited. Review findings are published in full through the \gls{gls:transparency-engine}. The Body may refer findings to the \gls{gls:constitutional-court} for constitutional review. Members are subject to the sector exclusion provisions of Article \ref{art:regulatory-bodies}, \ref{art:regulatory-bodies:s-sector-exclusion} -- persons with current or recent employment relationships with defense contractors or the military are excluded from the qualification pool.

\ArtSection{Military Budget Cap}{art:defense:s-military-budget-cap}
Military expenditure may not exceed two percent of Commonwealth GDP in any fiscal year. Expenditure above one and one-half percent requires a supermajority legislative vote with published justification. Any increase beyond two percent requires supermajority legislative vote, \gls{gls:constitutional-court} review confirming consistency with the defensive sufficiency principle, and a national referendum. The budget cap is cross-referenced with Article \ref{art:finance-money}'s fiscal provisions.

\ArtSection{Disarmament Commitment}{art:defense:s-disarmament-commitment}
The Commonwealth actively pursues verified multilateral nuclear disarmament and commits to reduce its arsenal proportionally to verified reductions by other nuclear states. This nuclear posture is a response to the current world, not a permanent preference.